\newpage
\chapter{KESIMPULAN DAN SARAN} \label{Bab V}

\section{Kesimpulan} \label{V.Kesimpulan}
Berdasarkan hasil penelitian yang telah dilaksanakan terhadap arsitektur 
\textit{ResNet-18} dan \textit{ResNet-50} dengan menerapkan augmentasi 
data kelas minoritas pada dataset DermaMNIST, dapat disimpulkan sebagai 
berikut.

\begin{enumerate}[noitemsep]
    \item Penerapan augmentasi data yang ditujukan khusus untuk kelas 
    minoritas memberikan dampak terhadap kedua performa model yaitu \textit{ResNet-18} 
    dan \textit{ResNet-50} dalam mengklasifikasikan data pada dataset 
    DermaMNIST. Kemampuan model meningkat dalam mempelajari kelas minoritas 
    karena didukung augmentasi sehingga proses pembelajaran menjadi 
    seimbang. Hal ini dapat dilihat dari hasil evaluasi pada model sendiri.

    \item Hasil menunjukkan performa yang berbeda terhadap \textit{ResNet-18} 
    dan \textit{ResNet-50} setelah melakukan proses \textit{training} 
    pada dataset DermaMNIST yang kelas minoritasnya sudah diterapkan 
    augmentasi data. Dari 4 skema yang telah dilakukan, model \textit{ResNet-18} 
    yang dilatih dengan citra 224x224 memiliki hasil yang lebih baik. 
    Model \textit{ResNet-18} dengan citra 224x224 mendapatkan akurasi 
    sebesar 0.80 (80\%), didukung dengan nilai \textit{precision} yang 
    mendapatkan nilai sebesar 0.71, nilai \textit{recall} mendapatkan 
    sebesar 0.76 dan nilai \textit{f1-score} sebesar 0.73. Selain itu, 
    model mampu dengan baik dalam memisahkan data dari masing-masing 
    kelasnya yang dilihat dari nilai \textit{AUC} sebesar 0.95. Nilai 
    MCC yang berada pada nilai 0.64 menunjukkan model mampu mengklasifikasi 
    dengan cukup baik dan tingkat korelasi positif yang kuat antara prediksi 
    dan aktual, serta model sudah tidak bergantung lagi denga kelas 
    mayoritas.

    \item Data uji digunakan untuk mengukur performa akhir model.
\end{enumerate}


% Berisi kesimpulan dari hasil dan pembahasan terkait penelitian yang dilakukan, dapat juga berupa temuan yang Anda dapatkan setelah melakukan penelitian atau analisis terhadap tugas akhir Anda. Memberikan jawaban dari poin pada subbab \nameref{I.Rumusan Masalah} dan \nameref{I.Tujuan}. \par

\section{Saran} \label{V.Saran}
Berisi saran mengenai aspek tugas akhir atau temuan yang dapat dikembangkan dan diperkaya di tugas akhir selanjutnya. Saran dapat berkaitan erat pada subbab \nameref{IV.Analisis}. \par